%
% kalkuel.tex
%
% (c) 2026 Prof Dr Andreas Müller
%
\section{Der Differentialkalkül}
Zur Zeit der Entstehung der Infinitesimalrechnung war die numerische
Berechnung meistens eine zu aufwendige Option.
An erster Stelle stand wann immer möglich die algebraische Rechnung.
Erst nach allen möglichen algebraischen Vereinfachungen war daran zu
denkten, konkrete Zahlenwerte einzusetzen und numerische Resultate
zu vergleichen.
Die Infinitesimalrechnung entstand daher zuerst primär als ein Kalkül,
was sich auch in dem im englischsprachigen Raum üblichen Namen
\emph{Calculus} widerspiegelt.

%
% Eulers Einführung in die Analysis
%
\subsection{Eulers Einführung in die Analysis}
Das Buch \emph{Introduction at analysin infinitorum}
\cite{buch:introductio1,buch:introductio2}
beginnt Leonhard Euler, den Leser in die Ideen der Inifinitesimalrechnung
einzuführen.
Es behandelt alle Aspekte der Infinitesimalrechnung, die ohne das
Konzept der Ableitung eingführt werden können.

Erst in seinen späteren Werk \cite{buch:institutiones}
\emph{Instituiones calculi differentialis} behandelt er die Ableitung
und das Integral.

Eulers Werk ist typische für die im 18. Jahrhunder vorherrschende
Betrachtungsweise der Infinitesimalrechnung.

%
% Die Rechenregeln der klassischen Infinitesimalrechnung
%
\subsection{Die Rechenregeln der klassischen Infinitesimalrechnung}

\begin{align}
\frac{d}{dx}(f(x)\pm g(x))
&=
\frac{d}{dx}f(x)
\pm
\frac{d}{dx}g(x)
\label{buch:einleitung:kalkuel:eqn:lin1}
\\
\frac{d}{dx}(\alpha f(x))
&=
\alpha \frac{d}{dx}f(x)
\label{buch:einleitung:kalkuel:eqn:lin2}
\\
\frac{d}{dx}\bigl(f(x)g(x)\bigr)
&=
\biggl(\frac{d}{dx}f(x)\biggr)
g(x)
+
f(x)
\biggl(\frac{d}{dx}g(x)\biggr)
\label{buch:einleitung:kalkuel:eqn:prod}
\\
\frac{d}{dx}c&=0\qquad\text{wenn $c$ eine Konstante ist}
\label{buch:einleitung:kalkuel:eqn:const}
\\
\frac{d}{dx}x&=1
\label{buch:einleitung:kalkuel:eqn:var}
\end{align}
Die Ableitung des reziproken $1/f(x)$ lässt sich daraus und aus
der Eigenschaft $f(x)\cdot 1/f(x)=1$ mit der Produktregel gewinnen,
ebenso die Quotientenregel.

In der klassischen Entwicklung der Infinitesimalrechnung werden diese
Regeln aus einre Definition der Ableitung als Grenzewert von
Differenzenquotienten gewonnen.
Man kann sich aber auch auf den Standpunkt stellen, dass diese Regeln
einen speziellen Operator auf allen Funktionen einer geeignet definierten
Funktionenmenge charakterisieren.
Die Funktionenmenge muss alle arithmetischen Operationen gestatten,
man muss also Funktionen addieren, subtrahieren, multiplizieren und
dividieren können.
Funktionen können mit Konstanten multipliziert werden.
Es muss möglich sein, Gleichungen mit Funktionsausdrücken aufzulösen.
Diese Eigenschaften lassen sich zusammenfassen in die Aussage, dass
die Funktionen ein Algebra über einem Zahlenkörper bilden müssen, wobei
man meistens von den reellen Zahlen als Zahlenkörper ausgeht.

Die Ableitung in einem Zahlenkörper ist dann ein Operator $D$, der auf
Funktionen wirkt.
Die beiden Regeln
\eqref{buch:einleitung:kalkuel:eqn:lin1}
und
\eqref{buch:einleitung:kalkuel:eqn:lin2}
erhalten die Form
\begin{align*}
 D(f+g) &= Df + Dg\\
 D(\alpha f) &\alpha Df,
\intertext{die besagt, dass $D$ ein linearer Operator ist.
Aus der Produktregel \eqref{buch:einleitung:kalkuel:eqn:prod}
wird}
D(fg) &= (Df)g+f(Dg).
\intertext{Diese Eigenschaft von $D$ heisst auch \emph{Derivation}.
Schliesslich identifiziert
\eqref{buch:einleitung:kalkuel:eqn:const}
die Konstanten und
\eqref{buch:einleitung:kalkuel:eqn:var}
wird zu}
Dx&=1,
\end{align*}
was charakterisiert die unabhängige Variable $x$, nach der der Operator
$D$ ableitet.
Die Vorgabe eines linearen Operators auf einer Algebra über den
rationalen Zahlen mit obigen Eigenschaften ist alles, was man braucht, um 
die Rechnungen der Infinitesimalrechnung durchführen zu können.
Sie entpuppt sich somit als hauptsächlich algebraischer Kalkül.



